% ===== §12 =====
\section{The Closure That Would Extinguish}
\label{p26:sec:deathtest}

\noindent
This section is the formal center of the paper. Its aim is to establish, by a thought experiment known to be counterfactual, that the closure of the gap would not fulfill drive and desire but extinguish them, and therefore that the non-closure of language is the formal condition of their persistence. The section states the counterfactual and its purpose, develops the consequence in three steps each given its own treatment, draws the reversal of the ordinary picture, and states the paper's central claim, that drive is the phenomenology of the limit.

\subsection*{The counterfactual and its use}

The section on the two faces established that complete symbolization is impossible. Suppose, against that, and only for the sake of the argument, that it were possible. Suppose the gap were closed, the remainder gathered into the signifier, the Real brought without loss into the symbolic. The supposition is counterfactual, and known to be so; the paper does not claim that the gap could be closed, and has argued that it cannot. The use of the supposition is different. By asking what closure would do, were it possible, the section shows what the non-closure secures, and thereby shows that the impossibility established earlier is not a deprivation but a condition of life. A thought experiment of this kind does not assert its premise. It reasons from a premise it rejects, in order to expose the value of the premise's denial.

\subsection*{First step, the drive circles because its object is absent}

The drive circles because its object is absent. This was established in the sections on the object and on the two movements: what the objet a marks is a void, and the drive turns around the void; the circling has its topology only because there is a want at the center and not a thing. The circuit is a circuit of an absence. Remove the absence and there is nothing to circle, no void to describe the loop around. The first step simply fixes what the circling depends on: not an object present at the center, but an object absent, a want given form. The drive's whole motion is predicated on the emptiness of the place it turns around.

\subsection*{Second step, closure fills the center}

Suppose now the gap closed. This is the counterfactual move, and its content is exact: the want at the center is filled, the absence made present, the void replaced by an attained object. Everything the earlier sections said could not happen is here supposed to happen, for the sake of seeing its consequence. The objet a, which was the form of a lack, would become a thing possessed; the remainder, which pressed from outside the signifier, would be gathered in; the Real, which was the internal limit of the symbolic, would be dissolved into the symbolic without residue. The center around which the drive turned would now be occupied. There would be, at the place of the want, a presence.

\subsection*{Third step, the circling stops and both movements die}

The drive no longer has a void to circle. Its circuit has lost the emptiness that gave it its shape, and a circuit around a filled center is no circuit at all; the loop collapses onto the point it enclosed. The circling therefore stops, and with it the drive is not fulfilled but extinguished. It is essential to see that this is extinction and not satisfaction. One might expect that to fill the want would be to satisfy the drive, since the drive seemed to press toward the want's filling; but the drive's satisfaction was in the circling, not in the filling, and to fill the center is to remove the condition of the circling and so to end the only satisfaction the drive had. Desire, for its part, has nowhere left to slide, since the chain of signifiers ran forward only because no signifier coincided with the thing, and now the coincidence has been supposed; the chain arrives, and arriving, halts. What the supposition delivers is not the satisfaction of drive and desire but their death. The state of total symbolization, imagined as a fulfillment, is in truth a stillness, a bond in which nothing further presses, nothing further turns, nothing further is generated.

\subsection*{The reversal}

The result reverses the ordinary picture. On the ordinary picture the gap is a defect and its closure would be a fulfillment, the moment the drive at last reached its object and desire at last came home. The thought experiment shows the opposite. The closure that looks like fulfillment is extinction, and the non-closure that looks like deprivation is the condition under which drive and desire go on at all. The impossibility of complete symbolization, established as impossibility in the section on the two faces, shows here its other aspect. It is not only that the gap cannot be closed. It is that its not being closed is what keeps the drive circling and the desire sliding, which is to say what keeps the relation in movement.

\begin{claim}[Closure would extinguish, not fulfill]
Were the gap closed, the drive would lose the void its circuit turns around and would stop, extinguished rather than satisfied, and desire would lose the non-coincidence that carries it along the chain and would halt. The non-closure of language is therefore not a deprivation but the formal condition of the persistence of drive and desire, and the impossibility of complete symbolization is the very thing that keeps a relation in movement.
\end{claim}

\begin{figure}[t!]
\centering
\begin{tikzpicture}[scale=0.95,
  lab/.style={font=\footnotesize\color{inksoft},align=center},
  drive/.style={{Latex[length=2mm]}-,gold,very thick}]
% LEFT: open, alive
\begin{scope}[shift={(0,0)}]
  \draw[gold,very thick,dashed] (0,0) circle (0.5cm);
  \node[font=\scriptsize\itshape\color{rosedeep}] at (0,0) {void};
  \draw[drive] (0,0) ++(40:1.15cm) arc (40:320:1.15cm);
  \node[lab] at (0,-2.1) {\textbf{open}\\ \tiny drive circles \\ \tiny \emph{alive}};
\end{scope}
% arrow of the counterfactual
\node[font=\footnotesize\itshape\color{rose},align=center] at (3.3,0) {suppose\\ closure\\ {\scriptsize(counterfactual)}};
\draw[-{Latex[length=2.5mm]},rose,thick] (1.8,0) -- (4.8,0);
% RIGHT: closed, extinguished
\begin{scope}[shift={(6.6,0)}]
  \fill[rose!35] (0,0) circle (0.5cm);
  \node[font=\scriptsize\color{inksoft}] at (0,0) {filled};
  \node[lab] at (0,-2.1) {\textbf{closed}\\ \tiny circuit collapses \\ \tiny \emph{extinguished}};
\end{scope}
\end{tikzpicture}
\caption{The counterfactual of closure. While the center is a void (left), the drive holds a circuit around it and the bond stays in movement. Were the gap filled (right), known to be impossible and supposed only for the argument, the circuit would have nothing to enclose and would collapse: closure delivers extinction, not fulfillment. The figure depicts what the non-closure secures by showing what its denial would cost.}
\label{p26:fig:closure}
\end{figure}

\subsection*{Drive as the phenomenology of the limit}

The central claim of the paper can now be stated. If the drive lives only where the gap is open, and dies where the gap is closed, then the drive is not a content that the gap withholds and a closure would deliver. The drive is the way the open gap shows itself, the felt form of the limit of language. It is not behind the limit, waiting; it is the limit, appearing.

\begin{claim}[Drive is the phenomenology of the limit]
Drive is not a content that language omits and might later recover. Drive is the phenomenology of manque, the felt form of the very limit of language. This is why the attempt to symbolize the gap returns not the drive but its renewed escape: what one tried to say was nothing other than what symbolization cannot close, so the saying finds the gap has moved, and the drive, which was that gap showing itself, has slipped once more beyond the reach of the word.
\end{claim}
